% 11_tweprints.tex — TwePrints: Communicating Research via Social Media Threads

\section{TwePrints: Communicating Research via Social Media Threads}
\label{sec:tweprints}

To facilitate early dissemination and improve research accessibility, an emerging practice among machine learning researchers is the use of \textit{TwePrints}---structured Twitter/X threads that summarize a paper's core contributions, often prior to formal publication.
TwePrints are not simply announcements; they aim to distill complex work into a concise and narrative-driven format, typically spanning 6--15 tweets.
Each post highlights a specific insight, result, or design decision, frequently accompanied by visuals, benchmarks, or interactive links to preprints, datasets, or code.

This medium serves several complementary functions.
First, it accelerates feedback cycles: by exposing work to the broader research community earlier, authors can identify ambiguities or receive informal peer commentary in advance of archival deadlines.
Second, TwePrints improve clarity and engagement by translating technical material into accessible formats for diverse audiences, including students, practitioners, and adjacent domain experts.
Third, they act as a durable reference artifact: summarizing experimental setups, dataset splits, agent behavior, or novel evaluation metrics in a format easily revisited and cited.

Successful examples of TwePrints include the thread accompanying \textit{VeriBench}, a benchmark that evaluates formal verification of AI-generated code using Lean~4.
Authored by the paper's contributors, this TwePrint distilled core findings---such as the sharp performance gains from feedback-driven agents, or the difficulty of verifying security-sensitive code---into a clear, narrative format.
The thread also clarified dataset design (across HumanEval, EasySet, CSSet, and SecuritySet), showcased the Trace agent's iterative loop, and highlighted relevant prior work.
To improve accessibility, the thread featured rich visuals, including result tables, example traces, and task illustrations.
By condensing complex material into a concise, scrollable sequence, the TwePrint facilitated early feedback, broadened visibility, and helped frame broader conversations around provable software generation and agentic code repair.
\textbf{WE CAN INCLUDE IMAGES FROM THE VERIBENCH TWEPRINT HERE.}

Well-constructed TwePrints follow an emerging structure: they begin with a high-level motivation, enumerate core contributions, visualize results or examples, and conclude with open questions or calls for replication.
While not a substitute for peer-reviewed publication, they are increasingly recognized as valuable complements---enhancing transparency, accelerating iteration, and widening the research conversation.
